\documentclass[11pt,a4paper]{article}
\usepackage[utf8]{inputenc}
\usepackage[T1]{fontenc}
\usepackage[spanish,es-minimal,es-noshorthands]{babel}
\usepackage{lmodern}
\usepackage[sfdefault,lining]{FiraSans}
\usepackage[a4paper,top=2.3cm,bottom=2.2cm,left=1.8cm,right=1.8cm,headheight=1.3cm,headsep=0.4cm,footskip=1cm]{geometry}
\usepackage{amsmath,amssymb}
\allowdisplaybreaks
\usepackage{icomma}
\usepackage{graphicx}
\usepackage[table]{xcolor}
\usepackage{tikz}
\usepackage{fontawesome5}
\usepackage{ccicons}
\usepackage{fancyhdr}
\usepackage{lastpage}
\usepackage{tasks}
\usepackage{needspace}
\usepackage[most]{tcolorbox}
\usepackage{array}
\usepackage[hidelinks]{hyperref}
\AtBeginDocument{\hypersetup{pdftitle={Soluciones de la ficha de la unidad 2: Sistemas de ecuaciones lineales. 2º Bach.},pdfsubject={Matemáticas CCSS 2º Bach.},pdfkeywords={sistemas de ecuaciones lineales, método de Gauss, método matricial, regla de Cramer, teorema de Rouché-Frobenius, sistemas homogéneos, discusión de sistemas, parámetros, problemas, PAU}}}

\definecolor{navy}{HTML}{1B2748}
\definecolor{accent}{HTML}{E8603A}
\definecolor{teal}{HTML}{1F9E90}
\definecolor{soft}{HTML}{EEF1F8}
\definecolor{muted}{HTML}{5D6685}

\setlength{\parindent}{0pt}
\setlength{\parskip}{3pt}
\settasks{label=\alph*), label-format=\bfseries\color{navy}, label-width=1.3em, item-indent=1.8em, column-sep=0.6em, after-item-skip=5pt, before-skip=3pt, after-skip=4pt}

% ---------- capçalera i peu ----------
\newcommand{\llicencia}{\href{https://creativecommons.org/licenses/by-nc-sa/4.0/deed.es}{\ccbyncsa}\ \ Licencia Creative Commons Atribución-NoComercial-CompartirIgual 4.0 Internacional (CC BY-NC-SA 4.0)}
\pagestyle{fancy}
\fancyhf{}
\fancyhead[L]{\raisebox{-0.35\height}{\includegraphics[height=0.95cm]{logo-gva.png}}}
\fancyhead[C]{\small\color{muted}Soluciones de la ficha · Unidad 2: Sistemas de ecuaciones lineales · 2\textsuperscript{o} Bach.}
\fancyhead[R]{\raisebox{-0.35\height}{\includegraphics[height=0.95cm]{logo-vermellar.png}}}
\fancyfoot[L]{\scriptsize\color{muted}\llicencia}
\fancyfoot[R]{\small\color{muted}\thepage}
\renewcommand{\headrulewidth}{0.4pt}
\renewcommand{\headrule}{\color{muted!40}\hrule height \headrulewidth}
\fancypagestyle{primera}{\fancyhf{}\renewcommand{\headrulewidth}{0pt}
  \fancyfoot[L]{\scriptsize\color{muted}\llicencia}
  \fancyfoot[R]{\small\color{muted}\thepage}}

% ---------- nivells i exercicis ----------
\newcommand{\estrelles}[1]{{\color{accent}\ifcase#1\or\faStar\faStar[regular]\faStar[regular]\or\faStar\faStar\faStar[regular]\or\faStar\faStar\faStar\fi}}
\newcounter{exer}
\newcommand{\ex}[1]{\par\needspace{4\baselineskip}\medskip\refstepcounter{exer}\noindent{\color{navy}\bfseries\theexer.}\ \ }

\newcommand{\seccio}[2]{\par\needspace{10\baselineskip}\bigskip
  \begin{tcolorbox}[colback=navy,colframe=navy,coltext=white,arc=2mm,boxrule=0pt,left=3mm,right=3mm,top=1.2mm,bottom=1.2mm]
  {\large\bfseries #1\quad #2}
  \end{tcolorbox}\vspace{-2pt}}


% quadre de recordatori (només a la fitxa d'exercicis)
\newcommand{\recorda}[1]{\par\smallskip\begin{tcolorbox}[colback=accent!7,colframe=accent!55,arc=2mm,boxrule=0.5pt,left=2.5mm,right=2.5mm,top=1mm,bottom=1mm]\small{\color{accent}\faLightbulb[regular]}\ \textbf{Recuerda:}\ #1\end{tcolorbox}}
% etiqueta dels exercicis de la PAU d'altres comunitats
\newcommand{\pau}[1]{{\color{teal}\small\bfseries\faUniversity\ PAU · #1.}\ }
\newcommand{\res}[1]{\boxed{#1}}\newcommand{\rest}[1]{\fbox{#1}}

% =====================================================================
% Capçalera FentMates de les fitxes: un color per a cada curs.
% Els colors són els de la pàgina inicial (LEVELS a assets/js/materials.js);
% el logo ve de _eines/logo/pdf/ (python3 _eines/logo/genera.py).
%
% Ús, al final del preàmbul de la unitat:
%   % =====================================================================
% Capçalera FentMates de les fitxes: un color per a cada curs.
% Els colors són els de la pàgina inicial (LEVELS a assets/js/materials.js);
% el logo ve de _eines/logo/pdf/ (python3 _eines/logo/genera.py).
%
% Ús, al final del preàmbul de la unitat:
%   % =====================================================================
% Capçalera FentMates de les fitxes: un color per a cada curs.
% Els colors són els de la pàgina inicial (LEVELS a assets/js/materials.js);
% el logo ve de _eines/logo/pdf/ (python3 _eines/logo/genera.py).
%
% Ús, al final del preàmbul de la unitat:
%   \input{../comu/fentmates.tex}
%   \fmcurs{1eso}
% i, al començament del document:
%   \capcalera{Matemàtiques}{1r ESO}{Fitxa d'exercicis · Unitat 2}{Divisibilitat}
% Opcional abans de \input: \def\fmneed{14} (espai mínim abans de cada secció).
% =====================================================================
\makeatletter
\def\fm@logodir{../../logo/pdf}
\providecommand{\fmneed}{10}
\def\fm@cursnom{}

% clau del curs, color i color del text damunt del color
\newcommand{\fm@curs}[3]{\@namedef{fm@color@#1}{#2}\@namedef{fm@text@#1}{#3}}
\fm@curs{1eso}{E6B84A}{navy}
\fm@curs{2eso}{3B82F6}{white}
\fm@curs{3eso}{FF7A50}{white}
\fm@curs{4eso-a}{A78BFA}{white}
\fm@curs{4eso-b}{6D28D9}{white}
\fm@curs{1batx-mi}{9B1C1C}{white}
\fm@curs{1batx-ccss}{EF6B6B}{white}
\fm@curs{2batx-mii}{1E6B47}{white}
\fm@curs{2batx-ccss}{2BB3A3}{white}

\newcommand{\fmcurs}[1]{%
  \@ifundefined{fm@color@#1}{\PackageError{fentmates}{Curs desconegut: #1}{}}{}%
  \edef\fm@tmp{\noexpand\definecolor{curs}{HTML}{\@nameuse{fm@color@#1}}%
               \noexpand\colorlet{curstext}{\@nameuse{fm@text@#1}}}\fm@tmp
  \colorlet{cursfosc}{curs!72!black}%
  \def\fm@logo{\fm@logodir/vertical-#1.pdf}%
  \def\fm@marca{\fm@logodir/marca-#1.pdf}}

% ---------- capçalera de la primera pàgina ----------
% #1 matèria (petit), #2 curs, #3 tipus de document i unitat, #4 títol
\newcommand{\capcalera}[4]{%
  \gdef\fm@cursnom{#2}%
  \par\noindent
  \begin{tcolorbox}[colback=curs,colframe=curs,boxrule=0pt,arc=2.2mm,width=3.5cm,height=2.15cm,
      valign=center,left=3mm,right=2mm,top=1mm,bottom=1mm,nobeforeafter,box align=center]
    \color{curstext}{\footnotesize #1}\par\vspace{1pt}{\fontsize{19}{21}\selectfont\bfseries #2\par}%
  \end{tcolorbox}\hspace{5mm}%
  \begin{minipage}[c]{\dimexpr\linewidth-3.5cm-5mm-2.9cm\relax}
    {\color{cursfosc}\small\bfseries\MakeUppercase{#3}}\par\vspace{2pt}
    {\color{navy}\fontsize{20}{23}\selectfont\bfseries #4\par}
  \end{minipage}\hfill
  \begin{minipage}[c]{2.6cm}\raggedleft\includegraphics[width=2.6cm]{\fm@logo}\end{minipage}\par
  \vspace{2pt}}

% ---------- barres de secció del color del curs ----------
\renewcommand{\seccio}[2]{\par\needspace{\fmneed\baselineskip}\bigskip
  \begin{tcolorbox}[colback=curs,colframe=curs,coltext=curstext,arc=2mm,boxrule=0pt,left=3mm,right=3mm,top=1.2mm,bottom=1.2mm]
  {\large\bfseries #1\quad #2}
  \end{tcolorbox}\vspace{-2pt}}

% ---------- estrelles de dificultat del color del curs (a Batxillerat no se'n posen) ----------
\@ifundefined{estrelles}{}{%
  \renewcommand{\estrelles}[1]{{\color{curs}\ifcase#1\or\faStar\faStar[regular]\faStar[regular]\or\faStar\faStar\faStar[regular]\or\faStar\faStar\faStar\fi}}}

% ---------- capçalera i peu de les pàgines ----------
% primera pàgina: logos institucionals al peu
\fancypagestyle{primera}{\fancyhf{}\renewcommand{\headrulewidth}{0pt}
  \fancyfoot[L]{\raisebox{-0.32\height}{\includegraphics[height=0.78cm]{logo-gva.png}}\hspace{5mm}%
    \raisebox{-0.32\height}{\includegraphics[height=0.78cm]{logo-vermellar.png}}}
  \fancyfoot[R]{\scriptsize\color{muted}\href{https://creativecommons.org/licenses/by-nc-sa/4.0/deed.ca}{\ccbyncsa}\ \ CC BY-NC-SA 4.0 · fentmates.org\qquad\small\thepage}}
% la resta: el símbol a l'esquerra i el curs a la dreta, amb el títol al centre com abans
\fancyhead[L]{\raisebox{-0.3\height}{\includegraphics[height=0.62cm]{\fm@marca}}}
\fancyhead[R]{\small\bfseries\color{cursfosc}\fm@cursnom}
\renewcommand{\headrule}{\color{curs!45}\hrule height \headrulewidth}
\makeatother

%   \fmcurs{1eso}
% i, al començament del document:
%   \capcalera{Matemàtiques}{1r ESO}{Fitxa d'exercicis · Unitat 2}{Divisibilitat}
% Opcional abans de \input: \def\fmneed{14} (espai mínim abans de cada secció).
% =====================================================================
\makeatletter
\def\fm@logodir{../../logo/pdf}
\providecommand{\fmneed}{10}
\def\fm@cursnom{}

% clau del curs, color i color del text damunt del color
\newcommand{\fm@curs}[3]{\@namedef{fm@color@#1}{#2}\@namedef{fm@text@#1}{#3}}
\fm@curs{1eso}{E6B84A}{navy}
\fm@curs{2eso}{3B82F6}{white}
\fm@curs{3eso}{FF7A50}{white}
\fm@curs{4eso-a}{A78BFA}{white}
\fm@curs{4eso-b}{6D28D9}{white}
\fm@curs{1batx-mi}{9B1C1C}{white}
\fm@curs{1batx-ccss}{EF6B6B}{white}
\fm@curs{2batx-mii}{1E6B47}{white}
\fm@curs{2batx-ccss}{2BB3A3}{white}

\newcommand{\fmcurs}[1]{%
  \@ifundefined{fm@color@#1}{\PackageError{fentmates}{Curs desconegut: #1}{}}{}%
  \edef\fm@tmp{\noexpand\definecolor{curs}{HTML}{\@nameuse{fm@color@#1}}%
               \noexpand\colorlet{curstext}{\@nameuse{fm@text@#1}}}\fm@tmp
  \colorlet{cursfosc}{curs!72!black}%
  \def\fm@logo{\fm@logodir/vertical-#1.pdf}%
  \def\fm@marca{\fm@logodir/marca-#1.pdf}}

% ---------- capçalera de la primera pàgina ----------
% #1 matèria (petit), #2 curs, #3 tipus de document i unitat, #4 títol
\newcommand{\capcalera}[4]{%
  \gdef\fm@cursnom{#2}%
  \par\noindent
  \begin{tcolorbox}[colback=curs,colframe=curs,boxrule=0pt,arc=2.2mm,width=3.5cm,height=2.15cm,
      valign=center,left=3mm,right=2mm,top=1mm,bottom=1mm,nobeforeafter,box align=center]
    \color{curstext}{\footnotesize #1}\par\vspace{1pt}{\fontsize{19}{21}\selectfont\bfseries #2\par}%
  \end{tcolorbox}\hspace{5mm}%
  \begin{minipage}[c]{\dimexpr\linewidth-3.5cm-5mm-2.9cm\relax}
    {\color{cursfosc}\small\bfseries\MakeUppercase{#3}}\par\vspace{2pt}
    {\color{navy}\fontsize{20}{23}\selectfont\bfseries #4\par}
  \end{minipage}\hfill
  \begin{minipage}[c]{2.6cm}\raggedleft\includegraphics[width=2.6cm]{\fm@logo}\end{minipage}\par
  \vspace{2pt}}

% ---------- barres de secció del color del curs ----------
\renewcommand{\seccio}[2]{\par\needspace{\fmneed\baselineskip}\bigskip
  \begin{tcolorbox}[colback=curs,colframe=curs,coltext=curstext,arc=2mm,boxrule=0pt,left=3mm,right=3mm,top=1.2mm,bottom=1.2mm]
  {\large\bfseries #1\quad #2}
  \end{tcolorbox}\vspace{-2pt}}

% ---------- estrelles de dificultat del color del curs (a Batxillerat no se'n posen) ----------
\@ifundefined{estrelles}{}{%
  \renewcommand{\estrelles}[1]{{\color{curs}\ifcase#1\or\faStar\faStar[regular]\faStar[regular]\or\faStar\faStar\faStar[regular]\or\faStar\faStar\faStar\fi}}}

% ---------- capçalera i peu de les pàgines ----------
% primera pàgina: logos institucionals al peu
\fancypagestyle{primera}{\fancyhf{}\renewcommand{\headrulewidth}{0pt}
  \fancyfoot[L]{\raisebox{-0.32\height}{\includegraphics[height=0.78cm]{logo-gva.png}}\hspace{5mm}%
    \raisebox{-0.32\height}{\includegraphics[height=0.78cm]{logo-vermellar.png}}}
  \fancyfoot[R]{\scriptsize\color{muted}\href{https://creativecommons.org/licenses/by-nc-sa/4.0/deed.ca}{\ccbyncsa}\ \ CC BY-NC-SA 4.0 · fentmates.org\qquad\small\thepage}}
% la resta: el símbol a l'esquerra i el curs a la dreta, amb el títol al centre com abans
\fancyhead[L]{\raisebox{-0.3\height}{\includegraphics[height=0.62cm]{\fm@marca}}}
\fancyhead[R]{\small\bfseries\color{cursfosc}\fm@cursnom}
\renewcommand{\headrule}{\color{curs!45}\hrule height \headrulewidth}
\makeatother

%   \fmcurs{1eso}
% i, al començament del document:
%   \capcalera{Matemàtiques}{1r ESO}{Fitxa d'exercicis · Unitat 2}{Divisibilitat}
% Opcional abans de \input: \def\fmneed{14} (espai mínim abans de cada secció).
% =====================================================================
\makeatletter
\def\fm@logodir{../../logo/pdf}
\providecommand{\fmneed}{10}
\def\fm@cursnom{}

% clau del curs, color i color del text damunt del color
\newcommand{\fm@curs}[3]{\@namedef{fm@color@#1}{#2}\@namedef{fm@text@#1}{#3}}
\fm@curs{1eso}{E6B84A}{navy}
\fm@curs{2eso}{3B82F6}{white}
\fm@curs{3eso}{FF7A50}{white}
\fm@curs{4eso-a}{A78BFA}{white}
\fm@curs{4eso-b}{6D28D9}{white}
\fm@curs{1batx-mi}{9B1C1C}{white}
\fm@curs{1batx-ccss}{EF6B6B}{white}
\fm@curs{2batx-mii}{1E6B47}{white}
\fm@curs{2batx-ccss}{2BB3A3}{white}

\newcommand{\fmcurs}[1]{%
  \@ifundefined{fm@color@#1}{\PackageError{fentmates}{Curs desconegut: #1}{}}{}%
  \edef\fm@tmp{\noexpand\definecolor{curs}{HTML}{\@nameuse{fm@color@#1}}%
               \noexpand\colorlet{curstext}{\@nameuse{fm@text@#1}}}\fm@tmp
  \colorlet{cursfosc}{curs!72!black}%
  \def\fm@logo{\fm@logodir/vertical-#1.pdf}%
  \def\fm@marca{\fm@logodir/marca-#1.pdf}}

% ---------- capçalera de la primera pàgina ----------
% #1 matèria (petit), #2 curs, #3 tipus de document i unitat, #4 títol
\newcommand{\capcalera}[4]{%
  \gdef\fm@cursnom{#2}%
  \par\noindent
  \begin{tcolorbox}[colback=curs,colframe=curs,boxrule=0pt,arc=2.2mm,width=3.5cm,height=2.15cm,
      valign=center,left=3mm,right=2mm,top=1mm,bottom=1mm,nobeforeafter,box align=center]
    \color{curstext}{\footnotesize #1}\par\vspace{1pt}{\fontsize{19}{21}\selectfont\bfseries #2\par}%
  \end{tcolorbox}\hspace{5mm}%
  \begin{minipage}[c]{\dimexpr\linewidth-3.5cm-5mm-2.9cm\relax}
    {\color{cursfosc}\small\bfseries\MakeUppercase{#3}}\par\vspace{2pt}
    {\color{navy}\fontsize{20}{23}\selectfont\bfseries #4\par}
  \end{minipage}\hfill
  \begin{minipage}[c]{2.6cm}\raggedleft\includegraphics[width=2.6cm]{\fm@logo}\end{minipage}\par
  \vspace{2pt}}

% ---------- barres de secció del color del curs ----------
\renewcommand{\seccio}[2]{\par\needspace{\fmneed\baselineskip}\bigskip
  \begin{tcolorbox}[colback=curs,colframe=curs,coltext=curstext,arc=2mm,boxrule=0pt,left=3mm,right=3mm,top=1.2mm,bottom=1.2mm]
  {\large\bfseries #1\quad #2}
  \end{tcolorbox}\vspace{-2pt}}

% ---------- estrelles de dificultat del color del curs (a Batxillerat no se'n posen) ----------
\@ifundefined{estrelles}{}{%
  \renewcommand{\estrelles}[1]{{\color{curs}\ifcase#1\or\faStar\faStar[regular]\faStar[regular]\or\faStar\faStar\faStar[regular]\or\faStar\faStar\faStar\fi}}}

% ---------- capçalera i peu de les pàgines ----------
% primera pàgina: logos institucionals al peu
\fancypagestyle{primera}{\fancyhf{}\renewcommand{\headrulewidth}{0pt}
  \fancyfoot[L]{\raisebox{-0.32\height}{\includegraphics[height=0.78cm]{logo-gva.png}}\hspace{5mm}%
    \raisebox{-0.32\height}{\includegraphics[height=0.78cm]{logo-vermellar.png}}}
  \fancyfoot[R]{\scriptsize\color{muted}\href{https://creativecommons.org/licenses/by-nc-sa/4.0/deed.ca}{\ccbyncsa}\ \ CC BY-NC-SA 4.0 · fentmates.org\qquad\small\thepage}}
% la resta: el símbol a l'esquerra i el curs a la dreta, amb el títol al centre com abans
\fancyhead[L]{\raisebox{-0.3\height}{\includegraphics[height=0.62cm]{\fm@marca}}}
\fancyhead[R]{\small\bfseries\color{cursfosc}\fm@cursnom}
\renewcommand{\headrule}{\color{curs!45}\hrule height \headrulewidth}
\makeatother

\fmcurs{2batx-ccss}

\begin{document}
\thispagestyle{primera}

\capcalera{Matemáticas CCSS}{2\textsuperscript{o} Bach.}{Soluciones de la ficha · Unidad 2}{Sistemas de ecuaciones lineales}
\vspace{8pt}

\begin{tcolorbox}[colback=soft,colframe=soft,arc=2mm,left=3mm,right=3mm,top=1.5mm,bottom=1.5mm]
\small Intenta resolver cada ejercicio antes de mirar su solución. Los pasos de Gauss pueden ser diferentes de los tuyos y ser igualmente correctos; lo que debe coincidir es la clasificación y la solución (en un SCI, la solución general se puede escribir de maneras diferentes según la incógnita que tomes como parámetro). Puedes comprobarlo con el resolutor de sistemas de los apuntes interactivos.
\end{tcolorbox}

% =====================================================
\seccio{1}{Sistemas de ecuaciones lineales. Método de Gauss}

\ex{1} \begin{tasks}(3)
  \task $5=5$, \ $-2=-2$, \ $7=7$: \rest{sí}
  \task La 2.\textsuperscript{a}: $2+1-2=1\neq-2$: \rest{no}
  \task La 1.\textsuperscript{a}: $0-1+3=2\neq5$: \rest{no}
\end{tasks}

\ex{1} \begin{tasks}(1)
  \task Sumando: $3x=3$, $x=1$, $y=2$. \rest{SCD}: rectas secantes en $(1,2)$.
  \task La 2.\textsuperscript{a} es la 1.\textsuperscript{a} por $-2$: \rest{SCI}, la misma recta: $x=1+2\lambda,\ y=\lambda$.
  \task La 1.\textsuperscript{a} por 2 da $6x+2y=4\neq5$: \rest{SI}, rectas paralelas.
\end{tasks}

\ex{2} Hay muchas respuestas posibles; por ejemplo:\begin{tasks}(1)
  \task La suma de las dos: \ $\res{3x+2z=5}$ \ (SCI).
  \task Mismos coeficientes que la 1.\textsuperscript{a} y otro término independiente: \ $\res{x+y+z=0}$ \ (SI).
  \task \ $\res{z=1}$: \ queda $x+y=2,\ 2x-y=1$, \ de donde $x=1,\ y=1,\ z=1$ (SCD).
\end{tasks}

\ex{1} \begin{tasks}(1)
  \task \[A^*=\left(\begin{array}{rrr|r}1 & 1 & 1 & 6\\2 & -1 & 3 & 9\\1 & 2 & -1 & 2\end{array}\right)\xrightarrow[F_3'=F_3-F_1]{F_2'=F_2-2F_1}\left(\begin{array}{rrr|r}1 & 1 & 1 & 6\\0 & -3 & 1 & -3\\0 & 1 & -2 & -4\end{array}\right)\]\[\xrightarrow{F_2\leftrightarrow F_3}\left(\begin{array}{rrr|r}1 & 1 & 1 & 6\\0 & 1 & -2 & -4\\0 & -3 & 1 & -3\end{array}\right)\]\[\xrightarrow{F_3''=F_3'+3F_2'}\left(\begin{array}{rrr|r}1 & 1 & 1 & 6\\0 & 1 & -2 & -4\\0 & 0 & -5 & -15\end{array}\right)\]$-5z=-15\Rightarrow z=3$; \ $y-6=-4\Rightarrow y=2$; \ $x=6-2-3$: \ $\res{x=1;\ y=2;\ z=3}$
  \task \[A^*=\left(\begin{array}{rrr|r}2 & 1 & -1 & 1\\1 & -1 & 2 & 5\\3 & 2 & 1 & 10\end{array}\right)\xrightarrow{F_1\leftrightarrow F_2}\left(\begin{array}{rrr|r}1 & -1 & 2 & 5\\2 & 1 & -1 & 1\\3 & 2 & 1 & 10\end{array}\right)\]\[\xrightarrow[F_3'=F_3-3F_1]{F_2'=F_2-2F_1}\left(\begin{array}{rrr|r}1 & -1 & 2 & 5\\0 & 3 & -5 & -9\\0 & 5 & -5 & -5\end{array}\right)\]\[\xrightarrow{F_3''=3F_3'-5F_2'}\left(\begin{array}{rrr|r}1 & -1 & 2 & 5\\0 & 3 & -5 & -9\\0 & 0 & 10 & 30\end{array}\right)\]$10z=30\Rightarrow z=3$; \ $3y-15=-9\Rightarrow y=2$; \ $x=5+2-6$: \ $\res{x=1;\ y=2;\ z=3}$
\end{tasks}

\ex{2} \begin{tasks}(1)
  \task \[A^*=\left(\begin{array}{rrr|r}1 & -1 & 2 & 1\\2 & 1 & -1 & 3\\4 & -1 & 3 & 5\end{array}\right)\xrightarrow[F_3'=F_3-4F_1]{F_2'=F_2-2F_1}\left(\begin{array}{rrr|r}1 & -1 & 2 & 1\\0 & 3 & -5 & 1\\0 & 3 & -5 & 1\end{array}\right)\]\[\xrightarrow{F_3''=F_3'-F_2'}\left(\begin{array}{rrr|r}1 & -1 & 2 & 1\\0 & 3 & -5 & 1\\0 & 0 & 0 & 0\end{array}\right)\]\rest{SCI}. Hacemos $z=\lambda$: $3y=1+5\lambda$, \ $x=1+y-2\lambda$: \ $\res{x=\tfrac{4-\lambda}{3};\ y=\tfrac{1+5\lambda}{3};\ z=\lambda}$
  \task \[A^*=\left(\begin{array}{rrr|r}1 & -1 & 2 & 1\\2 & 1 & -1 & 3\\4 & -1 & 3 & 6\end{array}\right)\xrightarrow[F_3'=F_3-4F_1]{F_2'=F_2-2F_1}\left(\begin{array}{rrr|r}1 & -1 & 2 & 1\\0 & 3 & -5 & 1\\0 & 3 & -5 & 2\end{array}\right)\]\[\xrightarrow{F_3''=F_3'-F_2'}\left(\begin{array}{rrr|r}1 & -1 & 2 & 1\\0 & 3 & -5 & 1\\0 & 0 & 0 & 1\end{array}\right)\]La última fila dice $0=1$: \ \rest{SI}, no tiene solución.
  \task La 2.\textsuperscript{a} y la 3.\textsuperscript{a} ecuaciones son la 1.\textsuperscript{a} multiplicada por $-2$ y por $3$: solo queda $x-2y+z=3$ (rango 1). \rest{SCI con dos parámetros}: \ $\res{x=3+2\lambda-\mu;\ y=\lambda;\ z=\mu}$
\end{tasks}

\ex{2} \pau{Navarra, EvAU extraordinaria 2023}\[A^*=\left(\begin{array}{rrr|r}1 & 1 & 3 & 1\\4 & 3 & 5 & 5\\2 & 1 & -1 & 3\end{array}\right)\xrightarrow[F_3'=F_3-2F_1]{F_2'=F_2-4F_1}\left(\begin{array}{rrr|r}1 & 1 & 3 & 1\\0 & -1 & -7 & 1\\0 & -1 & -7 & 1\end{array}\right)\]\[\xrightarrow{F_3''=F_3'-F_2'}\left(\begin{array}{rrr|r}1 & 1 & 3 & 1\\0 & -1 & -7 & 1\\0 & 0 & 0 & 0\end{array}\right)\]Queda una fila de ceros: \rest{SCI}. Con $z=\lambda$: \ $-y-7\lambda=1\Rightarrow y=-1-7\lambda$; \ $x=1-y-3\lambda=2+4\lambda$.\[\boxed{x=2+4\lambda;\ \ y=-1-7\lambda;\ \ z=\lambda,\quad\lambda\in\mathbb{R}}\]


% =====================================================
\seccio{2}{Expresión matricial. Método matricial}

\ex{2} \[\begin{pmatrix}2 & 0 & -1\\1 & 1 & 0\\0 & 1 & 1\end{pmatrix}\cdot\begin{pmatrix}x\\ y\\ z\end{pmatrix}=\begin{pmatrix}1\\ 0\\ 2\end{pmatrix}\]$|A|=\begin{vmatrix}2 & 0 & -1\\1 & 1 & 0\\0 & 1 & 1\end{vmatrix}=1\neq0$, por tanto $A$ tiene inversa:\[\mathrm{Adj}(A)=\begin{pmatrix}1 & -1 & 1\\-1 & 2 & -2\\1 & -1 & 2\end{pmatrix},\qquad A^{-1}=\frac{1}{1}\big(\mathrm{Adj}(A)\big)^t=\begin{pmatrix}1 & -1 & 1\\-1 & 2 & -1\\1 & -2 & 2\end{pmatrix}\]\[X=A^{-1}\cdot B=\begin{pmatrix}1 & -1 & 1\\-1 & 2 & -1\\1 & -2 & 2\end{pmatrix}\begin{pmatrix}1\\ 0\\ 2\end{pmatrix}=\begin{pmatrix}3\\ -3\\ 5\end{pmatrix}\ \Rightarrow\ \boxed{x=3;\ y=-3;\ z=5}\]

\ex{1} $\res{\left\{\begin{aligned}x+2y&=4\\-x+3z&=1\end{aligned}\right.}$ \ \ Tiene \rest{2 ecuaciones y 3 incógnitas} (el producto de una matriz $2\times3$ por una $3\times1$ es $2\times1$).\par No puede ser SCD: con 2 ecuaciones y 3 incógnitas, al escalonar siempre sobra al menos una incógnita. Las filas de $A$ no son proporcionales y, con $F_2'=F_2+F_1$, queda $2y+3z=5$: no aparece ninguna fila del tipo $0=d\neq0$. Por tanto, es \rest{SCI}, con un parámetro: \ $x=-1+3\lambda;\ y=\tfrac{5-3\lambda}{2};\ z=\lambda$.

\ex{2} \begin{tasks}(1)
  \task $|A|=1$, \ $A^{-1}=\begin{pmatrix}1 & -2\\-1 & 3\end{pmatrix}$, \ $X=\begin{pmatrix}1 & -2\\-1 & 3\end{pmatrix}\begin{pmatrix}7\\ 3\end{pmatrix}=\begin{pmatrix}1\\ 2\end{pmatrix}$: \ $\res{x=1;\ y=2}$
  \task $|A|=2$, \ $A^{-1}=\frac12\begin{pmatrix}1 & 1 & -1\\-1 & 1 & 1\\1 & -1 & 1\end{pmatrix}$, \ $X=\frac12\begin{pmatrix}1 & 1 & -1\\-1 & 1 & 1\\1 & -1 & 1\end{pmatrix}\begin{pmatrix}3\\ 3\\ 4\end{pmatrix}=\frac12\begin{pmatrix}2\\ 4\\ 4\end{pmatrix}$: \ $\res{x=1;\ y=2;\ z=2}$
\end{tasks}

\ex{2} \pau{Galicia, ABAU extraordinaria 2023}\[\begin{pmatrix}2 & 1 & 1\\0 & 2 & -1\\1 & 1 & 0\end{pmatrix}\cdot\begin{pmatrix}x\\ y\\ z\end{pmatrix}=\begin{pmatrix}2\\ -1\\ 1\end{pmatrix}\]$|A|=-1\neq0$, por tanto $A$ tiene inversa:\[\mathrm{Adj}(A)=\begin{pmatrix}1 & -1 & -2\\1 & -1 & -1\\-3 & 2 & 4\end{pmatrix},\qquad A^{-1}=\frac{1}{-1}\big(\mathrm{Adj}(A)\big)^t=\begin{pmatrix}-1 & -1 & 3\\1 & 1 & -2\\2 & 1 & -4\end{pmatrix}\]\[X=A^{-1}\cdot B=\begin{pmatrix}-1 & -1 & 3\\1 & 1 & -2\\2 & 1 & -4\end{pmatrix}\begin{pmatrix}2\\ -1\\ 1\end{pmatrix}=\begin{pmatrix}2\\ -1\\ -1\end{pmatrix}\ \Rightarrow\ \boxed{x=2;\ y=-1;\ z=-1}\]

\ex{2} \pau{País Vasco, PAU extraordinaria 2025}\begin{tasks}(1)
  \task $\begin{pmatrix}3 & -2 & -2\\1 & 0 & -1\\0 & 2 & -1\end{pmatrix}\cdot\begin{pmatrix}x\\ y\\ z\end{pmatrix}=\begin{pmatrix}3\\ 1\\ 0\end{pmatrix}$
  \task $|A|=3\cdot2+(-2)\cdot1+(-2)\cdot2=6-2-4=0$: \ \rest{no tiene inversa} y no podemos aplicar el método matricial.
  \task \[A^*=\left(\begin{array}{rrr|r}3 & -2 & -2 & 3\\1 & 0 & -1 & 1\\0 & 2 & -1 & 0\end{array}\right)\xrightarrow{F_1\leftrightarrow F_2}\left(\begin{array}{rrr|r}1 & 0 & -1 & 1\\3 & -2 & -2 & 3\\0 & 2 & -1 & 0\end{array}\right)\]\[\xrightarrow{F_2'=F_2-3F_1}\left(\begin{array}{rrr|r}1 & 0 & -1 & 1\\0 & -2 & 1 & 0\\0 & 2 & -1 & 0\end{array}\right)\]\[\xrightarrow{F_3'=F_3+F_2'}\left(\begin{array}{rrr|r}1 & 0 & -1 & 1\\0 & -2 & 1 & 0\\0 & 0 & 0 & 0\end{array}\right)\]\rest{SCI}. Con $z=\lambda$: \ $2y=\lambda$, \ $x=1+z$: \ $\res{x=1+\lambda;\ y=\tfrac{\lambda}{2};\ z=\lambda,\ \ \lambda\in\mathbb{R}}$
\end{tasks}

\ex{3} \begin{tasks}(1)
  \task $|A|=\begin{vmatrix}1 & k & 0\\0 & 1 & 1\\k & 0 & 1\end{vmatrix}=1+k^2>0$ para todo $k$: \ $A$ \rest{siempre es invertible}.
  \task $|A|=2$, \ $A^{-1}=\frac12\begin{pmatrix}1 & -1 & 1\\1 & 1 & -1\\-1 & 1 & 1\end{pmatrix}$, \ $X=\frac12\begin{pmatrix}1 & -1 & 1\\1 & 1 & -1\\-1 & 1 & 1\end{pmatrix}\begin{pmatrix}1\\ 2\\ 3\end{pmatrix}=\frac12\begin{pmatrix}2\\ 0\\ 4\end{pmatrix}$: \ $\res{x=1;\ y=0;\ z=2}$
\end{tasks}


% =====================================================
\seccio{3}{Regla de Cramer}

\ex{1} \begin{tasks}(1)
  \task $|A|=\begin{vmatrix}2 & -3\\5 & 1\end{vmatrix}=17$: \ $x=\dfrac{\begin{vmatrix}-4 & -3\\7 & 1\end{vmatrix}}{17}=\dfrac{17}{17}$, \ $y=\dfrac{\begin{vmatrix}2 & -4\\5 & 7\end{vmatrix}}{17}=\dfrac{34}{17}$: \ $\res{x=1;\ y=2}$
  \task $|A|=7$, \ $|A_x|=7$, \ $|A_y|=-7$, \ $|A_z|=14$: \ $\res{x=1;\ y=-1;\ z=2}$
\end{tasks}

\ex{2} \pau{Castilla y León, PAU ordinaria 2025}Queda $\left\{\begin{aligned}x+2y+z&=1\\x-2y+2z&=4\\3x+y-2z&=5\end{aligned}\right.$, con $|A|=\begin{vmatrix}1 & 2 & 1\\1 & -2 & 2\\3 & 1 & -2\end{vmatrix}=25\neq0$: es de Cramer.\[x=\frac{\begin{vmatrix}1 & 2 & 1\\4 & -2 & 2\\5 & 1 & -2\end{vmatrix}}{25}=\frac{52}{25},\quad y=\frac{\begin{vmatrix}1 & 1 & 1\\1 & 4 & 2\\3 & 5 & -2\end{vmatrix}}{25}=\frac{-17}{25},\quad z=\frac{\begin{vmatrix}1 & 2 & 1\\1 & -2 & 4\\3 & 1 & 5\end{vmatrix}}{25}=\frac{7}{25}\]\[\boxed{x=\tfrac{52}{25};\ \ y=-\tfrac{17}{25};\ \ z=\tfrac{7}{25}}\]

\ex{2} $|A|=\begin{vmatrix}1 & -1 & 2\\3 & 1 & -1\\2 & 3 & 1\end{vmatrix}=23\neq0$. \ $z=\dfrac{\begin{vmatrix}1 & -1 & 0\\3 & 1 & 2\\2 & 3 & 4\end{vmatrix}}{23}=\res{\dfrac{6}{23}}$; \ $x=\dfrac{10}{23}$, \ $y=\dfrac{22}{23}$ (y se comprueba en las tres ecuaciones).

\ex{3} $\begin{vmatrix}1&2\\ 2&-1\end{vmatrix}=-5\neq0$: \ $\left\{\begin{aligned}x+2y&=3+z\\ 2x-y&=1-3z\end{aligned}\right.$\[x=\frac{\begin{vmatrix}3+z&2\\ 1-3z&-1\end{vmatrix}}{-5}=\frac{-5+5z}{-5}=1-z,\qquad y=\frac{\begin{vmatrix}1&3+z\\ 2&1-3z\end{vmatrix}}{-5}=\frac{-5-5z}{-5}=1+z\]\[\boxed{x=1-\lambda;\ \ y=1+\lambda;\ \ z=\lambda,\quad\lambda\in\mathbb{R}}\]


% =====================================================
\seccio{4}{Sistemas homogéneos}

\ex{1} \begin{tasks}(1)
  \task $|A|=\begin{vmatrix}1 & 2 & -1\\2 & -1 & 1\\3 & 1 & 1\end{vmatrix}=-5\neq0$: \ \rest{SCD}, solo la solución trivial \ $\res{x=y=z=0}$.
  \task $|A|=0$ y $\begin{vmatrix}1&-1\\ 2&1\end{vmatrix}=3\neq0$: \ $\operatorname{rg}(A)=2<3$, \ \rest{SCI}. Sumando las dos primeras, $3x=0$; y de la 1.\textsuperscript{a}, $y=z$: \ $\res{x=0;\ y=\lambda;\ z=\lambda}$
\end{tasks}

\ex{2} \begin{tasks}(1)
  \task $|A|=\begin{vmatrix}1 & 1 & -2\\3 & -1 & 0\\-1 & 3 & -4\end{vmatrix}=0\Rightarrow\operatorname{rg}(A)<3$: \ \rest{SCI}, tiene infinitas soluciones además de la trivial.
  \task \[A^*=\left(\begin{array}{rrr|r}1 & 1 & -2 & 0\\3 & -1 & 0 & 0\\-1 & 3 & -4 & 0\end{array}\right)\xrightarrow[F_3'=F_3+F_1]{F_2'=F_2-3F_1}\left(\begin{array}{rrr|r}1 & 1 & -2 & 0\\0 & -4 & 6 & 0\\0 & 4 & -6 & 0\end{array}\right)\]\[\xrightarrow{F_3''=F_3'+F_2'}\left(\begin{array}{rrr|r}1 & 1 & -2 & 0\\0 & -4 & 6 & 0\\0 & 0 & 0 & 0\end{array}\right)\]Con $z=2\lambda$: \ $-4y+12\lambda=0\Rightarrow y=3\lambda$; \ $x=-y+2z=\lambda$: \ $\res{x=\lambda;\ y=3\lambda;\ z=2\lambda,\ \ \lambda\in\mathbb{R}}$
  \task $\lambda+3\lambda+2\lambda=6\lambda=12\Rightarrow\lambda=2$: \ $\res{x=2;\ y=6;\ z=4}$
\end{tasks}

\ex{1} Es homogéneo (compatible) y $\operatorname{rg}(A)\leq2<3$ incógnitas: \rest{SCI}. Pasando $z$ al segundo miembro: $2x-y=-z$, $x+3y=2z$ \ $\Rightarrow x=-\frac z7,\ y=\frac{5z}{7}$. Con $z=7\lambda$: \ $\res{x=-\lambda;\ y=5\lambda;\ z=7\lambda}$


% =====================================================
\seccio{5}{Teorema de Rouché-Frobenius}

\ex{1} \label{ex:rf}\begin{tasks}(1)
  \task $\begin{vmatrix}1&2\\ 3&-1\end{vmatrix}=-7\neq0\Rightarrow\operatorname{rg}(A)=2$; \ $|A^*|=\begin{vmatrix}1 & 2 & 5\\3 & -1 & 1\\2 & 1 & 4\end{vmatrix}=0\Rightarrow\operatorname{rg}(A^*)=2=n$: \ \rest{SCD} (la solución es $x=1,\ y=2$).
  \task $\operatorname{rg}(A)=2$; \ $|A^*|=\begin{vmatrix}1 & 2 & 5\\3 & -1 & 1\\1 & 1 & 2\end{vmatrix}=7\neq0\Rightarrow\operatorname{rg}(A^*)=3$: \ \rest{SI}.
  \task $|A|=0$ y $\begin{vmatrix}1&1\\ 2&-1\end{vmatrix}=-3\neq0\Rightarrow\operatorname{rg}(A)=2$. Los menores de orden 3 de $A^*$ con la columna de términos independientes también valen 0 $\Rightarrow\operatorname{rg}(A^*)=2<3$: \ \rest{SCI} (la 3.\textsuperscript{a} ecuación es la 1.\textsuperscript{a} más el doble de la 2.\textsuperscript{a}).
\end{tasks}

\ex{2} \pau{Murcia, PAU ordinaria 2025}$|A|=\begin{vmatrix}3 & 2 & -1\\1 & -1 & 2\\2 & 3 & -1\end{vmatrix}=-10\neq0\Rightarrow\operatorname{rg}(A)=\operatorname{rg}(A^*)=3=n$: \ \rest{SCD}. Por Cramer o por Gauss: \ $\res{x=1;\ y=1;\ z=2}$

\ex{2} $\begin{vmatrix}1 & 1 & 1\\1 & -1 & 1\\2 & 1 & -1\end{vmatrix}=6\neq0\Rightarrow\operatorname{rg}(A)=3$; \ $|A^*|=0\Rightarrow\operatorname{rg}(A^*)=3=n$: \ \rest{SCD}. Suprimimos la 4.\textsuperscript{a} ecuación (no forma parte del menor) y resolvemos las tres primeras: \ $\res{x=1;\ y=1;\ z=1}$

\ex{2} \begin{tasks}(1)
  \task \rest{Falsa}: si $\operatorname{rg}(A)=\operatorname{rg}(A^*)=2$ es SCD (como en el primer apartado del ejercicio \ref{ex:rf}).
  \task \rest{Verdadera}: $\operatorname{rg}(A)=3$ y $\operatorname{rg}(A^*)$ no puede pasar de 3 (tiene 3 filas): $\operatorname{rg}(A)=\operatorname{rg}(A^*)=3=n$.
  \task \rest{Falsa}: los rangos son distintos y el sistema es incompatible (ninguna solución).
  \task \rest{Falsa}: $\operatorname{rg}(A)\leq2<3$, así que, si es compatible, es indeterminado.
\end{tasks}


% =====================================================
\seccio{6}{Discusión de sistemas con parámetros}

\ex{2} \pau{Asturias, PAU extraordinaria 2025}\begin{tasks}(1)
  \task $|A|=\begin{vmatrix}1&-1\\ 10&10+m\end{vmatrix}=20+m$. \ Tiene solución única (SCD) si y solo si \ $\res{m\neq-20}$. (Para $m=-20$: $x-y=20$ y $x-y=15$, SI.)
  \task Debe ser $|A|=0$: $\res{m=-20}$; entonces la 2.\textsuperscript{a} ecuación es $x-y=15$ y, para que coincida con la 1.\textsuperscript{a}, $\res{\lambda=-15}$. Una solución: $\res{x=15,\ y=0}$ (en general, $x=15+t,\ y=t$).
\end{tasks}

\ex{2} \pau{Asturias, EBAU extraordinaria 2024}\begin{tasks}(1)
  \task $A\cdot B=\begin{pmatrix}1&2m+1\\ -1-m&0\end{pmatrix}$, \ $A\cdot B-C=\begin{pmatrix}1&1\\ m-1&1\end{pmatrix}$: \ \ $\res{\left\{\begin{aligned}x+y&=1\\ (m-1)x+y&=1\end{aligned}\right.}$
  \task $|A|=2-m$. \ Si $m\neq2$: \rest{SCD}, solución única. \ Si $m=2$: las dos ecuaciones son $x+y=1$, \rest{SCI}. \ Tiene solución \rest{para todo $m$}, pero solo es única si $m\neq2$. \ Para $m=-1$: $x+y=1$, $-2x+y=1$ \ $\Rightarrow\res{x=0;\ y=1}$.
\end{tasks}

\ex{2} \pau{Madrid, EvAU extraordinaria 2022}\begin{tasks}(1)
  \task $|A|=\begin{vmatrix}1&a&1\\ 1&0&-a\\ 1&1&1\end{vmatrix}=1-a^2=0\Rightarrow a=\pm1$.\par $\bullet$ $a\neq\pm1$: \rest{SCD}. \ $\bullet$ $a=1$: la 1.\textsuperscript{a} y la 3.\textsuperscript{a} ecuaciones son iguales, $\operatorname{rg}(A)=\operatorname{rg}(A^*)=2$: \rest{SCI}. \ $\bullet$ $a=-1$: $\operatorname{rg}(A)=2$ y $\begin{vmatrix}1&-1&2\\ 1&0&0\\ 1&1&2\end{vmatrix}=4\neq0\Rightarrow\operatorname{rg}(A^*)=3$: \rest{SI}.
  \task $x+z=2$, $x=0$, $x+y+z=2$: \ $\res{x=0;\ y=0;\ z=2}$
\end{tasks}

\ex{2} \pau{Madrid, EvAU extraordinaria 2024}\begin{tasks}(1)
  \task $|A|=7a-21=0\Rightarrow a=3$. \ $\bullet$ $a\neq3$: \rest{SCD}. \ $\bullet$ $a=3$: $\begin{vmatrix}2&1\\ 3&1\end{vmatrix}=-1\neq0$, $\operatorname{rg}(A)=2$, y los menores de orden 3 de $A^*$ valen 0 (la 3.\textsuperscript{a} ecuación es $E_1+2E_2$): $\operatorname{rg}(A^*)=2$, \rest{SCI}.
  \task Restando $E_2-E_1$: $x-2z=-2$. Con $z=\lambda$: $x=-2+2\lambda$, \ $y=-3x-z=6-7\lambda$:\par $\res{x=-2+2\lambda;\ y=6-7\lambda;\ z=\lambda}$
\end{tasks}

\ex{2} \pau{Madrid, EvAU ordinaria 2024}\begin{tasks}(1)
  \task $|A|=2a^2-2a=2a(a-1)=0\Rightarrow a=0,\ a=1$.\par $\bullet$ $a\neq0,1$: \rest{SCD}. \ $\bullet$ $a=0$: $\operatorname{rg}(A)=2$ y $\begin{vmatrix}2&1&0\\ 1&0&1\\ 1&1&2\end{vmatrix}=-3\neq0$: \rest{SI}. \ $\bullet$ $a=1$: la 2.\textsuperscript{a} y la 3.\textsuperscript{a} ecuaciones coinciden ($x+y+z=2$), $\operatorname{rg}(A)=\operatorname{rg}(A^*)=2$: \rest{SCI}.
  \task $2x+y+z=1$ y $x+y+z=2$. Restando: $x=-1$; con $z=\lambda$: \ $\res{x=-1;\ y=3-\lambda;\ z=\lambda}$
\end{tasks}

\ex{2} \pau{Madrid, PAU extraordinaria 2025}\begin{tasks}(1)
  \task $|A|=2-a=0\Rightarrow a=2$. \ $\bullet$ $a\neq2$: \rest{SCD}. \ $\bullet$ $a=2$: la 3.\textsuperscript{a} ecuación, $3x+2z=5$, es la suma de las dos primeras: $\operatorname{rg}(A)=\operatorname{rg}(A^*)=2$, \rest{SCI}.
  \task Con $z=\lambda$: $x+y=2-\lambda$, $2x-y=3-\lambda$ $\Rightarrow$ \ $\res{x=\tfrac{5-2\lambda}{3};\ y=\tfrac{1-\lambda}{3};\ z=\lambda}$
\end{tasks}

\ex{2} \pau{Murcia, PAU extraordinaria 2025}\begin{tasks}(1)
  \task $|A|=2-a=0\Rightarrow a=2$. \ $\bullet$ $a\neq2$: \rest{SCD}. \ $\bullet$ $a=2$: $\operatorname{rg}(A)=2$ y $\begin{vmatrix}1&-1&1\\ 2&-1&2\\ 3&2&4\end{vmatrix}=1\neq0\Rightarrow\operatorname{rg}(A^*)=3$: \rest{SI}.
  \task $E_2-E_1$: $x=0$. Entonces $-y+z=1$ y $2y-z=4$: \ $\res{x=0;\ y=5;\ z=6}$
\end{tasks}

\ex{2} \pau{Murcia, EBAU ordinaria 2023}$|A|=-a^2+3a-2=-(a-1)(a-2)=0\Rightarrow a=1,\ a=2$.\par $\bullet$ $a\neq1,2$: \rest{SCD}. \ $\bullet$ $a=1$: la 1.\textsuperscript{a} y la 3.\textsuperscript{a} ecuaciones tienen los mismos coeficientes y términos independientes distintos: \rest{SI}. \ $\bullet$ $a=2$: la 1.\textsuperscript{a} y la 2.\textsuperscript{a} son iguales, $\operatorname{rg}(A)=\operatorname{rg}(A^*)=2$: \rest{SCI}.\par Para $a=0$: $y+z=1$, $2x+y+z=1$, $x+z=2$ \ $\Rightarrow\res{x=0;\ y=-1;\ z=2}$

\ex{2} $|A|=\begin{vmatrix}1&1&-1\\ 2&-1&k\\ 1&2&1\end{vmatrix}=-k-8=0\Rightarrow\res{k=-8}$. Si $k\neq-8$: \rest{SCD}, solo la solución trivial. Para $k=-8$ es \rest{SCI}: de la 1.\textsuperscript{a}, $z=x+y$; en la 3.\textsuperscript{a}, $2x+3y=0$.\[\boxed{x=3\lambda;\ \ y=-2\lambda;\ \ z=\lambda,\quad\lambda\in\mathbb{R}}\]

\ex{3} \pau{Madrid, EvAU extraordinaria 2023}\begin{tasks}(1)
  \task $|A|=a^3-7a+6=(a-1)(a-2)(a+3)=0\Rightarrow a=1,\ a=2,\ a=-3$.\par $\bullet$ $a\neq1,2,-3$: \rest{SCD}. \ $\bullet$ $a=1$: la 1.\textsuperscript{a} y la 2.\textsuperscript{a} ecuaciones son iguales, $\operatorname{rg}(A)=\operatorname{rg}(A^*)=2$: \rest{SCI}. \ $\bullet$ $a=2$: la 2.\textsuperscript{a} y la 3.\textsuperscript{a} tienen los mismos coeficientes y términos independientes distintos: \rest{SI}. \ $\bullet$ $a=-3$: $\operatorname{rg}(A)=2$ y $\begin{vmatrix}-3&1&1\\ 1&-3&-3\\ 1&2&1\end{vmatrix}=-8\neq0$: \rest{SI}.
  \task $y+2z=1$, $x+2z=0$, $x+2y=1$ \ $\Rightarrow\res{x=-\tfrac13;\ y=\tfrac23;\ z=\tfrac16}$
\end{tasks}

\ex{3} \pau{Castilla y León, PAU ordinaria 2025}En general, $|A|=ab^2-2a-5b+13$ (el ejercicio de la sección de Cramer resolvía el caso $a=1$, $b=-2$).\begin{tasks}(1)
  \task $|A|=18-a=0\Rightarrow a=18$. \ $a\neq18$: \rest{SCD}. \ $a=18$: $\begin{vmatrix}1&-1\\ 3&1\end{vmatrix}=4\neq0$, $\operatorname{rg}(A)=2$, y $\begin{vmatrix}2&1&1\\ -1&2&4\\ 1&-1&5\end{vmatrix}=36\neq0$: \rest{SI}.
  \task $|A|=3b^2-5b+7$, que no se anula nunca (discriminante $25-84<0$): \rest{SCD para todo $b$}.
\end{tasks}


% =====================================================
\seccio{7}{Problemas de planteamiento}

\ex{1} \pau{Canarias, EBAU ordinaria 2022}$x$, $y$, $z$: personas de pádel surf, kayak y moto. \ \[\left\{\begin{aligned}x+y+z&=45\\40x+20y+60z&=1700\\x-3y&=0\end{aligned}\right.\] \ Con $x=3y$: $4y+z=45$ y $7y+3z=85$\par $\Rightarrow$ \ \rest{30 de pádel surf, 10 de kayak y 5 de moto acuática}.

\ex{2} \pau{Canarias, EBAU ordinaria 2023}$x$: niños y jóvenes, $y$: adultos, $z$: jubilados. \ $\frac y4=\frac{2x}{5}$, \ $\frac{z}{x+y}=\frac{9}{26}$: \ \[\left\{\begin{aligned}x+y+z&=875\\-8x+5y&=0\\9x+9y-26z&=0\end{aligned}\right.\] \ De la 3.\textsuperscript{a} y la 1.\textsuperscript{a}: $35z=9\cdot875$, $z=225$; $x+y=650$ e $y=\frac85x$\par $\Rightarrow$ \ \rest{250 niños y jóvenes, 400 adultos y 225 jubilados}.

\ex{1} \pau{Canarias, EBAU extraordinaria 2024}\[\left\{\begin{aligned}300x+400y+800z&=28000\\x-2y&=0\\y-2z&=0\end{aligned}\right.\] ($x$ móviles, $y$ tabletas, $z$ portátiles). \ Con $y=2z$, $x=4z$: $2800z=28\,000$\par $\Rightarrow$ \ \rest{40 móviles, 20 tabletas y 10 portátiles}.

\ex{2} \pau{Canarias, EBAU ordinaria 2024}$x$, $y$, $z$: torneos locales, nacionales e internacionales. \ \[\left\{\begin{aligned}x-2y&=0\\2y-5z&=0\\0{,}3x+0{,}3y+0{,}3z&=51\end{aligned}\right.\] \ El total es $51/0{,}3=170$; con $x=2y$ y $z=\frac25y$: $\frac{17}{5}y=170$\par $\Rightarrow$ \ \rest{100 locales, 50 nacionales y 20 internacionales}.

\ex{2} \pau{Cantabria, EBAU extraordinaria 2022}$x$, $y$, $z$: precio de un helado, de un granizado y de una horchata. Le han cobrado 4 helados, 2 granizados y 3 horchatas por 20 €; lo que tenía que pagar es 15 €; y el 50\,\% de $x+y$ es 2 €:\ \[\left\{\begin{aligned}4x+2y+3z&=20\\3x+2y+2z&=15\\x+y&=4\end{aligned}\right.\] \ Restando las dos primeras: $x+z=5$; con $y=4-x$, $z=5-x$ en la 2.\textsuperscript{a}: $x=3$\par $\Rightarrow$ \ \rest{helado 3 €, granizado 1 €, horchata 2 €}.

\ex{1} \pau{Cantabria, EBAU ordinaria 2023}Un saco de 25 kg vale 100 €, uno de 50 kg, 200 € y uno de 100 kg, 400 €: \ \[\left\{\begin{aligned}x+y+z&=180\\100x+200y+400z&=29200\\x-2y-2z&=0\end{aligned}\right.\] \ De la 1.\textsuperscript{a} y la 3.\textsuperscript{a}: $3(y+z)=180$, $y+z=60$, $x=120$; e $y+2z=86$\par $\Rightarrow$ \ \rest{120 sacos de 25 kg, 34 de 50 kg y 26 de 100 kg}.

\ex{2} \pau{Cantabria, EBAU extraordinaria 2024}\begin{tasks}(1)
  \task $x$ banderas, $y$ camisetas, $z$ gorras: \ \[\left\{\begin{aligned}5x+6y+2z&=500\\2y-z&=0\\x+y&=70\end{aligned}\right.\]
  \task $|A|=\begin{vmatrix}5 & 6 & 2\\0 & 2 & -1\\1 & 1 & 0\end{vmatrix}=-5\neq0$: $\operatorname{rg}(A)=\operatorname{rg}(A^*)=3$, \rest{SCD}.
  \task \rest{40 banderas, 30 camisetas y 60 gorras}.
\end{tasks}

\ex{2} \pau{Cantabria, EBAU ordinaria 2024}\begin{tasks}(1)
  \task $x$, $y$, $z$: €/kg de tierra, piedras y semillas en A. El proveedor B cobra $13\,000-8\,800=4\,200$ €: \ \[\left\{\begin{aligned}300x+200y+100z&=13000\\100x+100y+20z&=4200\\2x+2y-z&=0\end{aligned}\right.\]
  \task $|A|=\begin{vmatrix}300 & 200 & 100\\100 & 100 & 20\\2 & 2 & -1\end{vmatrix}=-14000\neq0$: \rest{SCD}.
  \task Sustituyendo $z=2x+2y$: $5x+4y=130$ y $x+y=30$\par $\Rightarrow$ \ \rest{tierra 10 €/kg, piedras 20 €/kg, semillas 60 €/kg}.
\end{tasks}

\ex{1} \pau{Castilla y León, EBAU ordinaria 2023}$x$ teatro, $y$ baloncesto, $z$ concierto: \ \[\left\{\begin{aligned}0{,}9x+0{,}9y+0{,}9z&=117\\2x-z&=0\\-y+z&=20\end{aligned}\right.\] \ $x+y+z=130$; con $x=\frac z2$, $y=z-20$: $\frac52z=150$\par $\Rightarrow$ \ \rest{teatro 30 €, baloncesto 40 €, concierto 60 €}.

\ex{1} \pau{Castilla-La Mancha, EvAU extraordinaria 2022}$x$ Isabel, $y$ Carmen, $z$ Enma: $x+1=3z$, $z=\frac34y$: \ \[\left\{\begin{aligned}x+y+z&=15\\x-3z&=-1\\-3y+4z&=0\end{aligned}\right.\]\par $\Rightarrow$ \ \rest{Isabel 8, Carmen 4 y Enma 3}.

\ex{2} \pau{Castilla-La Mancha, EvAU extraordinaria 2022}\[\left\{\begin{aligned}10x+6y+3z&=851000\\x-y&=2000\\5x-6z&=13000\end{aligned}\right.\] \ Con $y=x-2000$ y $z=\frac{5x-13000}{6}$ en la 1.\textsuperscript{a}: $18{,}5x=869\,500$.\par $\Rightarrow$ \ \rest{deportivo 47\,000 €, familiar 45\,000 €, monovolumen 37\,000 €}.

\ex{1} \pau{Castilla-La Mancha, EvAU ordinaria 2022}\[\left\{\begin{aligned}12x+6y+6z&=51\\x-y&=1\\y-z&=-0{,}5\end{aligned}\right.\] \ Con $x=y+1$, $z=y+0{,}5$: $24y+15=51$\par $\Rightarrow$ \ \rest{negro 2{,}50 €, con leche 1{,}50 €, blanco 2 €}.

\ex{2} \pau{Castilla-La Mancha, EvAU extraordinaria 2023}\[\left\{\begin{aligned}x+y+z&=50\\2x-y&=1\\2x-5z&=0\end{aligned}\right.\] \ Con $y=2x-1$, $z=\frac{2x}{5}$: $\frac{17}{5}x=51$\par $\Rightarrow$ \ \rest{15 de blanco, 29 de rosado y 6 de tinto}.

\ex{2} \pau{Castilla-La Mancha, EvAU ordinaria 2023}$z-y=3(y-x)$: \ \[\left\{\begin{aligned}x+y+z&=70000\\x+y-z&=0\\3x-4y+z&=0\end{aligned}\right.\] \ Sumando las dos primeras, $z=35\,000$\par $\Rightarrow$ \ \rest{15\,000 del I, 20\,000 del II y 35\,000 del III}.

\ex{2} \pau{Castilla-La Mancha, EvAU extraordinaria 2024}$x$ Alba, $y$ Pablo, $z$ David: \ \[\left\{\begin{aligned}x+y+z&=72\\-x+y-z&=0\\x-\tfrac{1}{3}y-z&=0\end{aligned}\right.\] \ Sumando las dos primeras: $2y=72$, $y=36$\par $\Rightarrow$ \ \rest{24 Alba, 36 Pablo y 12 David}.

\ex{2} \pau{Castilla-La Mancha, EvAU ordinaria 2024}$x$ oro, $y$ plata, $z$ bronce: $z+2=2(y-2)$: \ \[\left\{\begin{aligned}x+y+z&=36\\3x-z&=0\\2y-z&=6\end{aligned}\right.\]\par $\Rightarrow$ \ \rest{6 de oro, 12 de plata y 18 de bronce}.

\ex{3} \pau{Castilla-La Mancha, EvAU ordinaria 2024}$x$ Inés, $y$ Nerea, $z$ nulos: $4y-3x=z+1$, \ $x=y+7$, \ $z=0{,}05(x+y+z)\Leftrightarrow x+y-19z=0$: \ \[\left\{\begin{aligned}-3x+4y-z&=1\\x-y&=7\\x+y-19z&=0\end{aligned}\right.\]\par $\Rightarrow$ \ \rest{Inés 32 votos, Nerea 25 votos y 3 nulos}.

\ex{1} \pau{Castilla-La Mancha, PAU extraordinaria 2025}\[\left\{\begin{aligned}x+y+z&=25\\140x+130y+160z&=3610\\x-4y+z&=0\end{aligned}\right.\] \ De la 1.\textsuperscript{a} y la 3.\textsuperscript{a}: $5y=25$, $y=5$; después $7x+8z=148$ y $x+z=20$\par $\Rightarrow$ \ \rest{12 italianos, 5 portugueses y 8 japoneses}.

\ex{2} \pau{Cataluña, PAU extraordinaria 2022}\begin{tasks}(1)
  \task $x$ autocar, $y$ a pie, $z$ taxi. El trayecto total es $x+y+z$: \ $x=\frac{x+y+z}{2}$, \ $y=\frac{x+y+z}{20}$, \ $x=z+5$:\ \[\left\{\begin{aligned}x-y-z&=0\\-x+19y-z&=0\\x-z&=5\end{aligned}\right.\] \ De la 1.\textsuperscript{a} y la 2.\textsuperscript{a}, $x=10y$ y $z=9y$; en la 3.\textsuperscript{a}, $y=5$: \ \rest{50 km en autocar, 5 km a pie y 45 km en taxi}.
  \task $\frac{50}{100}+\frac55+\frac{45}{90}=0{,}5+1+0{,}5=\res{2\text{ horas}}$.
\end{tasks}

\ex{3} \pau{Cataluña, PAU ordinaria 2023}\begin{tasks}(1)
  \task $x$ termómetro, $y$ mascarilla, $z$ test: \ \[\left\{\begin{aligned}5x+30y+10z&=550\\x+15y+20z&=200\end{aligned}\right.\] \ Dos ecuaciones y tres incógnitas, $\operatorname{rg}(A)=\operatorname{rg}(A^*)=2<3$: \rest{SCI}, \textbf{no} se pueden determinar los precios. Restando $5E_2-E_1$: $45y+90z=450$: \ $\res{y=10-2z;\ \ x=50+10z}$.
  \task $z=4$: \ \rest{la mascarilla cuesta 2 € y el termómetro, 90 €}.
\end{tasks}

\ex{1} \pau{Madrid, EvAU extraordinaria 2024}$x$ guitarras, $y$ pianos, $z$ violines: \ \[\left\{\begin{aligned}x+y+z&=70\\x-y-z&=0\\4x+2y+z&=180\end{aligned}\right.\] \ De las dos primeras, $x=35$\par $\Rightarrow$ \ \rest{35 guitarras, 5 pianos y 30 violines}.

\ex{1} \pau{Madrid, EvAU ordinaria 2024}\[\left\{\begin{aligned}x+y+z&=600\\10x+8y+5z&=4900\\x-2z&=0\end{aligned}\right.\] \ Con $x=2z$: $3z+y=600$ y $25z+8y=4900$\par $\Rightarrow$ \ \rest{200 generales, 300 de estudiante y 100 infantiles}.

\ex{2} \pau{Navarra, PAU ordinaria 2025}Las cajas valen 6 €, 12 € y 24 €: \ \[\left\{\begin{aligned}x+y+z&=65\\x-y&=5\\6x+12y+24z&=720\end{aligned}\right.\]\par $\Rightarrow$ \ \rest{30 cajas de 250 g, 25 de 500 g y 10 de 1 kg} (números naturales: la solución tiene sentido).

\ex{1} \pau{País Vasco, PAU extraordinaria 2025}\[\left\{\begin{aligned}x+y+z&=8\\10x+20y+50z&=290\\x-2y&=0\end{aligned}\right.\] \ Con $x=2y$: $3y+z=8$ y $40y+50z=290$\par $\Rightarrow$ \ \rest{2 billetes de 10 €, 1 de 20 € y 5 de 50 €}.

\ex{2} \pau{País Vasco, PAU ordinaria 2025}\[\left\{\begin{aligned}x+y+z&=20\\0{,}25x+0{,}5y+z&=10\\x+1{,}8y+3{,}3z&=35{,}6\end{aligned}\right.\] \ Multiplicando por 4 la 2.\textsuperscript{a} y por 10 la 3.\textsuperscript{a}, y restando la 1.\textsuperscript{a}: $y+3z=20$ y $8y+23z=156$\par $\Rightarrow$ \ \rest{8 latas de A, 8 de B y 4 de C}.

\ex{3} \pau{Asturias, PAU ordinaria 2025}\begin{tasks}(1)
  \task $x+3y=90$, $2x+15y=225$\par $\Rightarrow$ \ \rest{VIP 75 €, normal 5 €}.
  \task Para $m\neq6$ es SCD con $x=\frac{45m}{m-6}$, $y=\frac{15(m-12)}{m-6}$. Si $6<m<12$, el precio $y$ es negativo: por ejemplo, $\res{m=9}$ da $x=135$, $y=-15$, que no tiene sentido.
  \task $|A|=m-6=0\Rightarrow\res{m=6}$: $2x+6y=90$ contradice $x+3y=90$ (que multiplicada por 2 da 180): \rest{SI}.
\end{tasks}

\ex{3} \pau{Asturias, EBAU extraordinaria 2023}\begin{tasks}(1)
  \task $\res{\left\{\begin{aligned}x-my&=0\\ 0{,}15x+0{,}15y&=3m\end{aligned}\right.}$
  \task $|A|=0{,}15(1+m)$. Si $m\neq-1$: \rest{SCD} (solución única, $x=\frac{20m^2}{1+m}$, $y=\frac{20m}{1+m}$). Si $m=-1$: $x+y=0$ y $x+y=-20$: \rest{SI}. \ Con 9 litros, $3m=9$, $m=3$: $x+y=60$, $x=3y$\par $\Rightarrow$ \ \rest{45 copas de vino tinto y 15 de blanco}.
\end{tasks}


\end{document}
